% Separate candidate only; NO automatic manuscript insertion.
\subsection{Lengthscale and the limits of amplitude transfer}
For fixed-smoothness Mat\'ern kernels, equivalent RKHS norms preserve a managed
spectral order that already holds; they do not identify its exponent or amplitude
from the kernel alone. If $\mu_j(\ell)\sim c(\ell)j^{-b}$ with common $b>1$,
the asymptotic minimizer of the theoretical envelope satisfies
\[
 \frac{\lambda_T^{\rm bound}(\ell_2)}{\lambda_T^{\rm bound}(\ell_1)}
 =\left(\frac{c(\ell_2)}{c(\ell_1)}\right)^{1/(b+1)}
  \left(\frac{\rho(\ell_2)}{\rho(\ell_1)}\right)^{b/(b+1)},\qquad\rho=B/A.
\]
This is not an oracle-penalty identity. Managed eigenvalues alone do not determine
expected-payoff signal or score noise, and two-sided spectral bounds need not supply
a limiting amplitude. In the real JKP cache, the initial Mat\'ern finite-range slopes
vary across lengthscales and the predeclared amplitude-identification gate fails.
Mechanical transfer paths are therefore falsification exercises. The 492 common
OOS months for decisions 1984--2024 include fixed lambda, annual CV, joint CV and
one-anchor controls; adaptive selection uses only the previous 60 genuinely OOS
candidate returns. Actual signed weights reconstruct net costs and paired circular
blocks assess differences conditional on fitted paths. The standalone investigation
does not establish a convincing structural one-constant contribution for the current
paper. We recommend exclusion, while retaining the negative result and its audit.
The Gaussian extension tests the same finite-range heuristic while preserving
the negative control against an assumed polynomial population law.
